\documentclass{article}
\usepackage{amsmath, amssymb, amsthm}
\usepackage{hyperref}
\usepackage{geometry}
\geometry{margin=1in}

\title{Erd\H{o}s Problem \#1199}
\date{Last updated: 2026-04-04}
\author{Source: erdosproblems.com}

\begin{document}
\maketitle

\noindent\textbf{Status:} OPEN \quad \textbf{No prize}

\noindent\textbf{Tags:} additive combinatorics, ramsey theory

\medskip
\noindent\textbf{Problem Statement:}

Is it true that in any $2$-colouring of $\mathbb{N}$ there exists an infinite set $A$ such that all elements of $A+A$ are the same colour?

\bigskip
\noindent\textbf{Remarks:}

A conjecture of Owings \cite{Ow74}. Hindman \cite{Hi79} has shown that this is false for $3$-colourings. 

If we do not ask for the elements $2a$ for $a\in A$ to also be the same colour the answer is yes by Hindman's theorem (see [532]).

\bigskip
\noindent\textbf{References:}

[Hi79] Hindman, Neil, Partitions and sums of integers with repetition. J. Combin. Theory Ser. A (1979), 19--32.
 
 [Ow74] J. Owings, E2494. Amer. Math. Monthly (1974), 902.

\bigskip
\noindent\small{Source: \url{https://www.erdosproblems.com/1199}}
\end{document}
